\documentclass[10pt]{letter}

\usepackage[margin=2.2cm]{geometry}
\usepackage[T1]{fontenc}
\usepackage{lmodern}
\usepackage{microtype}
\usepackage[colorlinks=true,allcolors=blue]{hyperref}

\signature{Corresponding author: [NAME]\\
Affiliation: [INSTITUTION]\\
Address: [POSTAL ADDRESS]\\
Email: [EMAIL]}
\address{[AUTHOR ACTION REQUIRED: corresponding author's address]}

\begin{document}
\setlength{\parskip}{0.25em}

\begin{letter}{Editors\\Nature Machine Intelligence}
\opening{Dear Editors,}

Please consider our Article, ``Token identity provides a fixed routing signal for residual experts in language models'', for publication in \emph{Nature Machine Intelligence}.

Mixture-of-experts language models generally learn a context-dependent router, making it difficult to separate conditional capacity from the routing decision. We present a controlled alternative in which every token retains a shared contextual MLP while token identity selects only a narrow residual path. In a parameter- and token-matched single-seed comparison at 306.5 million parameters and 8 billion training tokens, fixed token routing finishes with lower negative log-likelihood on a fixed diagnostic stream drawn from the training split than dense computation. Final checkpoints remain near parity on three zero-shot tasks. Routing has lower WikiText-2 perplexity and C4-validation NLL, while dense has lower NLL on a Pile test subset; the effect is therefore corpus dependent. The routed implementation also trains more slowly. A separate learned contextual router leads dense at short budget but trails its matched dense control at 1 billion tokens.

We believe this result will interest the journal's broad machine-learning readership because it separates two usually coupled components of MoE systems: conditional capacity and adaptive routing. The finding does not claim that fixed routing is universally superior; it shows that stable token identity can allocate useful residual capacity when a shared path preserves contextual computation. The manuscript makes the non-held-out diagnostic stream, mixed independent-corpus results, single-seed design and efficiency limits explicit and includes reproducible source artifacts.

The manuscript is submitted for double-anonymized peer review. Author names, affiliations and contribution statements are listed below and have been omitted from the reviewer manuscript.

\textbf{Authors and affiliations:} [AUTHOR ACTION REQUIRED: list every author, primary affiliation, current address if applicable, and ORCID.]\par
\textbf{Author contributions:} [AUTHOR ACTION REQUIRED: provide a CRediT-style statement.]\par
\textbf{Competing interests:} [AUTHOR ACTION REQUIRED: confirm that none exist or provide the complete declaration.]\par
\textbf{Related manuscripts:} [AUTHOR ACTION REQUIRED: confirm that the overlapping TMLR submission has been withdrawn or finally decided before Nature submission, and disclose any related manuscript under consideration or in press.]\par
\textbf{Prior editorial discussions:} [AUTHOR ACTION REQUIRED: state whether any discussion with a Nature Machine Intelligence editor occurred.]

Generative artificial-intelligence tools assisted language editing, consistency checks against author-provided CSV/JSON metric files, and figure and revision scripting; the authors verified all results, numerical values, scientific claims and final text, take full responsibility, and disclose this use in Methods.

Thank you for considering our manuscript.

\closing{Sincerely,}
\end{letter}

\end{document}
